\begin{abstract}
Iterative code repair---using feedback signals to guide LLM re-generation on failing problems---improves functional correctness, but which feedback signal to use remains an open question. We compare execution test feedback against pylint/mypy static analysis feedback for Llama 3.1 8B Instruct in an iso-compute setting ($B$=1000 output tokens per problem) on HumanEval and MBPP. Execution feedback significantly outperforms pylint/mypy on both benchmarks (McNemar's test: HumanEval $p$=0.0001, MBPP $p$<10$^{-18}$), improving MBPP pass@1 by +40.2pp while pylint/mypy repair actively \emph{reduces} HumanEval pass@1 by $-$4.3pp. The mechanism is a style-function dissociation: pylint flags 100\% of HumanEval failures, but 94.3\% of those flags are Convention-category style rules (missing newlines, missing docstrings) that fire regardless of functional correctness, while functional Error/Warning coverage is only 12.5\%. On complex algorithmic tasks, style-guided rewrites corrupt logical structure; on simpler function-completion tasks (MBPP), they provide marginal benefit. These findings suggest that feedback signal selection for LLM code repair should prioritize functional informativeness over aggregate coverage.
\end{abstract}
